%% Guideline 5a: Chapter 1 must be the Introduction and be labelled so.
\chapter{Introduction}
\label{ch:introduction}

This chapter introduces the topic. It gives the background to the problem so that readers
understand why the work matters, relates the thesis to other work in the field, and points
to the main findings and contributions.

Write paragraphs as normal text with a blank line between them. The template adds the space
between paragraphs and the one and a half line spacing for you, so never add manual spacing.

\section{Problem Statement}
State clearly the problem your work answers. Keep it focused and concise.

\section{Objectives}
List specific, measurable and achievable objectives, for example:
\begin{enumerate}
  \item Build a test set for the task.
  \item Compare two or more methods on that test set.
  \item Analyse where each method succeeds and fails.
\end{enumerate}

\section{Thesis Organisation}
Give a short description of each chapter. Chapter~\ref{ch:literature} reviews related work,
Chapter~\ref{ch:methodology} describes the method, Chapter~\ref{ch:experiments} the
experiments, Chapter~\ref{ch:results} the results and Chapter~\ref{ch:conclusions} the
conclusions.

\section{Summary}
Each chapter may end with a short summary (optional).
